%%%%%%%%%%%%%%%%%%%%%%%%%%%%%%%%%%%%%%%%%%%%%%%%%%%%%%%%%%%%%%%%%%%%%%%%

%%% LaTeX Template for AAMAS-2026 (based on sample-sigconf.tex)
%%% Prepared by the AAMAS-2026 Publication Chairs based on the version from AAMAS-2025. 

%%%%%%%%%%%%%%%%%%%%%%%%%%%%%%%%%%%%%%%%%%%%%%%%%%%%%%%%%%%%%%%%%%%%%%%%

\documentclass[sigconf,anonymous]{aamas} 

\usepackage{balance} 
\usepackage{algorithm}
\usepackage{algorithmic}
\usepackage{booktabs}
\usepackage{multirow}
\usepackage{colortbl}
\usepackage{xcolor}
\usepackage[section]{placeins}
\usepackage{amsmath}
\DeclareMathOperator*{\argmax}{arg\,max}
\DeclareMathOperator*{\argmin}{arg\,min}
\usepackage{caption}
\usepackage{subcaption}
\usepackage{todonotes}
\usepackage{orcidlink}
\usepackage{hyperref}


%%%%%%%%%%%%%%%%%%%%%%%%%%%%%%%%%%%%%%%%%%%%%%%%%%%%%%%%%%%%%%%%%%%%%%%%

\setcopyright{ifaamas}
\acmConference[AAMAS '26]{Proc.\@ of the 25th International Conference
on Autonomous Agents and Multiagent Systems (AAMAS 2026)}{May 25 -- 29, 2026}
{Paphos, Cyprus}{C.~Amato, L.~Dennis, V.~Mascardi, J.~Thangarajah (eds.)}
\copyrightyear{2026}
\acmYear{2026}
\acmDOI{}
\acmPrice{}
\acmISBN{}


%%%%%%%%%%%%%%%%%%%%%%%%%%%%%%%%%%%%%%%%%%%%%%%%%%%%%%%%%%%%%%%%%%%%%%%%

\acmSubmissionID{<<submission id>>}

\title[Counterfactual Gradient Alignment]{Counterfactual Gradient Alignment: Optimizing Directional Expert Supervision for Data-Efficient Learning}

\author{Arthur Pendragon}
\affiliation{
  \institution{Camelot Castle}
  \city{Camelot}
  \country{United Kingdom}}
\email{king.arthur@camelot.uk}

\author{Nimue}
\affiliation{
  \institution{The Lady's Lake}
  \city{Avalon}
  \country{United Kingdom}}
\email{lady.of.the.lake@avalon.uk}

\begin{abstract}
Typically, machine learning algorithms must find patterns in samples drawn from an unknown data distribution, with no prior contextual information about the input space or intuition about the problem they are solving. Humans, by contrast, can apply both intuition and prior knowledge to quickly solve new problems, requiring significantly less training data than modern deep learning architectures. In this paper we investigate how to address this disparity by enriching traditional labels for supervised classification tasks; embedding additional information through gradient-based \emph{expert annotations}. We present a framework for \textbf{Counterfactual Gradient Alignment (CGA)} which (1) annotates existing observations with directions in the feature space pointing towards proximal counterfactuals and (2) utilizes a novel family of loss functions---including ReLU, Softplus, and Sign-based variants---which reward model gradients that align with these directions. We demonstrate the effectiveness of CGA across a diverse suite of benchmarks: synthetic 2D datasets (XOR, Circles, TwoMoons), image classification (MNIST), and multiple natural language sentiment classification tasks (IMDB, SST-2, SST-5). Our results reveal significant accuracy improvements in low-data regimes, such as a +7.17\% gain on IMDB and a +14.5\% improvement on complex 2D boundaries. Furthermore, we uncover a critical synergy between CGA and active learning, showing that directional supervision is uniquely effective when applied to informative boundary samples, whereas it can act as noise on random data subsets. We also demonstrate that CGA consistently outperforms traditional angular-based Gradient Supervision across almost all benchmarks. The proposed method provides a blueprint for more collaborative and data-efficient human-AI teaching loops.
\end{abstract}

\keywords{Counterfactual Explanations, Gradient Supervision, Active Learning, Data Efficiency, Human-AI Interaction}

%%%%%%%%%%%%%%%%%%%%%%%%%%%%%%%%%%%%%%%%%%%%%%%%%%%%%%%%%%%%%%%%%%%%%%%%

\newcommand{\BibTeX}{\rm B\kern-.05em{\sc i\kern-.025em b}\kern-.08em\TeX}

%%%%%%%%%%%%%%%%%%%%%%%%%%%%%%%%%%%%%%%%%%%%%%%%%%%%%%%%%%%%%%%%%%%%%%%%

\begin{document}

\pagestyle{fancy}
\fancyhead{}

\maketitle 

%%%%%%%%%%%%%%%%%%%%%%%%%%%%%%%%%%%%%%%%%%%%%%%%%%%%%%%%%%%%%%%%%%%%%%%%

\section{Introduction}
Typically, machine learning algorithms must find patterns in samples drawn from an unknown data distribution, with no prior contextual information about the input space or intuition about the problem they are solving. Humans, by contrast, can apply both intuition and prior knowledge to quickly solve new problems, requiring significantly less training data than modern deep learning architectures. In this paper we investigate how to address this disparity by enriching traditional labels for supervised classification tasks; embedding additional information through gradient-based \emph{expert annotations}. We present a framework for \textbf{Counterfactual Gradient Alignment (CGA)} which (1) annotates existing observations with directions in the feature space pointing towards proximal counterfactuals and (2) utilizes a novel family of loss functions---including ReLU, Softplus, and Sign-based variants---which reward model gradients that align with these directions. We demonstrate the effectiveness of CGA across a diverse suite of benchmarks: synthetic 2D datasets (XOR, Circles, TwoMoons), image classification (MNIST), and multiple natural language sentiment classification tasks (IMDB, SST-2, SST-5). Our results reveal significant accuracy improvements in low-data regimes, such as a +7.17\% gain on IMDB and a +14.5\% improvement on complex 2D boundaries. Furthermore, we uncover a critical synergy between CGA and active learning, showing that directional supervision is uniquely effective when applied to informative boundary samples, whereas it can act as noise on random data subsets. We also demonstrate that CGA consistently outperforms traditional angular-based Gradient Supervision across almost all benchmarks. The proposed method provides a blueprint for more collaborative and data-efficient human-AI teaching loops.

\subsection{The Bandwidth of Supervision Signal}
In the traditional supervised learning paradigm, the supervisory signal is remarkably sparse. For each training instance, the model receives only a single categorical label. This label serves as a discrete marker in the input space but conveys almost nothing about the local topology of the decision boundary around that point. The model must essentially "guess" the orientation and curvature of the decision surface by aggregating many of these sparse labels over time. This process is inherently inefficient, particularly in high-dimensional manifolds where the number of possible decision surfaces that satisfy the discrete labels is vast.

We propose that closing the data efficiency gap requires increasing the bandwidth of the supervision signal between the expert and the model. Instead of providing only a label, we should enable experts to communicate their structural and geometric intuition directly to the model. An expert who labels a review as "Positive" sentiment already understands the linguistic geometry of the sentence. They know which emotive words are pivotal and how their negation or replacement would "flip" the sentiment to "Negative." If this directional knowledge could be captured in a differentiable form, it could provide a much denser training signal, regularizing the model's derivative landscape and guiding it toward a more semantically plausible decision boundary from the very first training samples.

\subsection{Counterfactual Gradient Alignment (CGA)}
In this work, we introduce \textbf{Counterfactual Gradient Alignment (CGA)}, a framework that converts directional expert intuition into machine-interpretable, first-order supervision. In the CGA framework, an expert identifies a direction vector $\mathbf{d}$ in the feature space pointing from a query instance to a proximal counterfactual instance---a point where the class membership changes. We then optimize the model's gradients to align with this vector, ensuring that moving in the specified direction consistently reduces the model's confidence in the original class. This process effectively "teaches" the model the local geometry of the decision boundary at each training point.

Where gradients have been used previously as a natural method for providing human-interpretable explanations of model behavior \cite{simonyan2013deep}, CGA aims to reverse this pipeline. We use expert-provided "explanations" (in the form of counterfactual directions) to provide direct supervision to the model's gradients. We evaluate CGA across a diverse benchmark suite encompassing vision, natural language processing, and synthetic manifolds. We introduce a family of optimized loss functions: ReLU, Softplus, and Sign-based variants. We compare our approach to the standard angular-based Gradient Supervision (GS) method proposed by \cite{teney2020learning} and demonstrate its superiority across almost all tasks.

\subsection{Synergy with Active Learning}
A primary discovery of our study is the profound synergy between CGA and active learning. We demonstrate that directional supervision is a high-precision signal that thrives when targeted at informative boundary samples identified via entropy-based importance sampling. Boundary samples are those for which the model is most uncertain about its local topology. Conversely, we show that applying directional alignment to random or "easy" samples can be counter-productive. Directional signals on central samples can act as noise, conflicting with the primary classification objective. This suggests that CGA is best deployed in a targeted, interactive "teaching" loop where the model proactively queries an expert for directional intuition on its most ambiguous boundaries. Our results show accuracy gains of up to +7.17\% on IMDB sentiment analysis and +14.5\% on complex 2D synthetic boundaries.

\section{Related Work}
\label{related_work}

The objective of integrating expert knowledge into machine learning models has led to several distinct research paradigms. Here, we review state-of-the-art methods related to human priors, counterfactuals, and gradient supervision.

\subsection{Embedding Human Priors}
Integrating human domain knowledge into model training is a longstanding goal in the field. \citeauthor{rieger2020interpretations} \cite{rieger2020interpretations} developed contextual decomposition explanation penalisation (CDEP), which enables embedding domain knowledge during training. CDEP works by penalising model explanations that do not align with those of an oracle, thereby forcing the model to ignore spurious correlations. Ross et al. similarly embedded human priors as explanations and develop a loss function which balances cross entropy loss with a penalty for misaligned feature attributions \cite{ross2017right}. These methods often focus on reinforcing the magnitude of gradients to align with pre-supposed importance weights. CGA extends this paradigm by focusing on the \emph{direction of class transitions}, providing a more direct geometric constraint on the model's decision surface.

\subsection{Counterfactually Augmented Data}
Kaushik et al. \cite{Kaushik2020Learning} established the utility of counterfactually augmented data (CAD) by tasking human annotators to minimally edit individual observations to produce a counterfactual. They showed that incorporating counterfactuals within the training data leads to improved generalisation and reduces sensitivity to spurious signals. In a follow-up paper, Kaushik et al. investigated why CAD reduces spurious correlations \cite{kaushik2021explaining} and identified human annotations as superior to automated augmentations. While their work treats counterfactuals as additional factual training points, CGA treats them as \emph{directional supervisors}. We utilize the \emph{relationship} between the original and counterfactual sample to supervise the model's derivative landscape directly. This provides a more continuous and structurally aware regularizing signal than point-based methods.

\subsection{Gradient Supervision Paradigms}
Teney et al. \cite{teney2020learning} introduce "gradient supervision", which involves penalising the angle between model gradients and counterfactual vectors. This approach aims to align the model's decision surface between pairs of counterfactual examples. Our work expands on this paradigm by evaluating magnitude-aware and bounded alignment signals (Softplus, Sign). We find these variants to be more robust across diverse dataset topologies and high-dimensional manifolds. Specifically, we move beyond pure angular alignment to consider the \emph{directional derivative}. This allows for more stable regularisation in high-noise environments where pure angular signals can be unstable.

\subsection{Interpretability as Supervision}
Recent research has explored using model interpretability techniques as a source of supervision. Simonyan et al. \cite{simonyan2013deep} used gradients to generate saliency maps for human interpretation of model focus. CGA reverses this relationship, using human-provided saliency (in the form of counterfactual directions) to supervise the model's gradients. This creates a bi-directional interface where human intuition and machine optimization can align more effectively. This is particularly important in low-data regimes where the categorical signal is insufficient to guide the model toward a semantically meaningful decision boundary.

\section{Gradient-Based Learning}
\label{sec:gradient_learning}

\subsection{Directional Expert Annotation}
In traditional supervised machine learning, we assume a dataset of $N$ random independent identically distributed (i.i.d.) observations {$\mathbf{x_i}, y_i$} $\sim P(\mathbf{x},y)$, where $\mathbf{x_i}\in\mathcal{R}^d$ and $y_i$ is the target class label \cite{hastie2009elements}. In the CGA framework, we consider triplets {$\mathbf{x_i}, y_i, \mathbf{d_i}$}, where $\mathbf{d_i} \in \mathcal{R}^d$ defines a unit vector from $\mathbf{x_i}$ pointing towards a proximal counterfactual $\hat{\mathbf{x}}_i$: 
\begin{equation}
    \mathbf{d_i}=\frac{\hat{\mathbf{x}}_i - \mathbf{x_i}}{||\hat{\mathbf{x}}_i - \mathbf{x_i}||}. \label{eq:dir_unit}
\end{equation}
We propose that moving from $\mathbf{x_i}$ towards $\hat{\mathbf{x}}_i$ should result in a decrease in the model's prediction probability for the true class $y$. Let $f_y(\mathbf{x})$ be the model's logit score for class $y$. We define the \textbf{Directional Derivative} $d$ as the projection of the gradient onto the counterfactual direction:
\begin{equation}
    d = \nabla_x f_{y}(x) \cdot \mathbf{d}_{i}
\end{equation}
This derivative quantifies how much the class confidence changes as we perturb the input along the expert-provided direction.

\subsection{Optimized Loss Function Variants}
We evaluate three loss variants $\mathcal{L}_{d}$ designed to minimize $d$, aligning the model's confidence drop with the counterfactual path. The intuition is that $d$ should be strongly negative at the counterfactual direction.

\paragraph{ReLU Directional Loss}
The simplest alignment objective is the ReLU loss:
\begin{equation}
    \mathcal{L}_{\text{ReLU}} = \max(0, d)
\end{equation}
This applies a linear penalty only if the model's confidence is incorrectly \emph{increasing} toward the counterfactual boundary. It is a "hinge-like" objective that stops regularizing once the gradient sign is correct. While computationally efficient, it may lack the signal necessary to encourage robust decision boundaries in extremely sparse data regimes.

\paragraph{Softplus Directional Loss}
To address the "zero-penalty" correct regions of ReLU, we propose the Softplus loss:
\begin{equation}
    \mathcal{L}_{\text{Softplus}} = \log(1 + e^{\beta d})
\end{equation}
Softplus is a smooth version of ReLU that provides a continuous alignment signal. Unlike ReLU, it provides a non-zero gradient even when $d$ is negative. This encourages the model to maintain a "steep" confidence drop-off even when it is already locally "correct". This property proved essential in high-dimensional text embedding spaces where decision boundaries are subtle.

\paragraph{Sign (Tanh) Directional Loss}
To ensure stability and prevent gradient explosion from outliers, we propose the Sign variant:
\begin{equation}
    \mathcal{L}_{\text{Sign}} = \tanh(\beta d) + 1
\end{equation}
This variant squashes the derivative magnitude, focusing purely on the sign of the alignment. It ensures a bounded, stable regularizing signal across different manifolds. The $\beta$ parameter acts as a temperature scale, controlling the "sharpness" of the directional signal.

\subsection{Total Training Objective}
The total loss is a weighted mixture of classification and alignment objectives:
\begin{align}\label{eq:lambda}
    \mathcal{L}_{\text{Total}} = (1-\alpha)\mathcal{L}_{\textrm{CE}} + \alpha \mathcal{L}_{d}
\end{align}
where $\alpha \in [0,1]$ controls the trade-off between standard classification and geometric supervision. In our JAX-based implementation, we leverage efficient automatic differentiation to compute the directional derivative $d$ within the optimization loop.

\section{Experimental Setup}
\label{sec:experimental_setup}

\subsection{Implementation and Optimization}
All experiments were implemented in the \textbf{JAX/Flax} ecosystem. This choice allowed us to leverage efficient automatic differentiation and vectorization (\texttt{vmap}). Gradient computation for directional derivatives was performed using JAX's \texttt{jacobian} function, allowing for precise first-order supervision without manual derivative definitions. We utilized the \textbf{AdamW} optimizer with weight decay $10^{-4}$ and $\beta_1=0.9, \beta_2=0.95$. The learning rate was adapted per-task, ranging from $10^{-4}$ for vision to $10^{-2}$ for synthetic tasks. All models were trained for 5-10 epochs to simulate rapid adaptation scenarios common in few-shot learning.

\subsection{Multimodal Benchmark Suite}
We evaluate CGA on three distinct benchmark categories: synthetic 2D manifolds, image classification, and natural language sentiment.

\paragraph{Synthetic 2D Topologies:}
We evaluate on XOR (four clusters), Concentric Circles (one class inside another), and Two Moons (interleaved arcs). We utilized an "optimum classifier inversion" to generate counterfactuals. The model is a simple 2-layer MLP with 64 units and tanh activation. These tasks allow us to justify design decisions by observing geometric alignment directly in the feature space.

\paragraph{Image Classification (MNIST):}
We use the standard MNIST digits dataset to evaluate how the quality of directional intuition affects CGA performance. We compare three counterfactual path generation strategies:
\begin{itemize}
    \item \textbf{FACE}: Plausible paths generated using the Feasible Path algorithm \cite{poyiadzi2020face}, ensuring transitions respect local data density.
    \item \textbf{PCA}: Shortest paths in PCA-transformed feature space.
    \item \textbf{Raw}: Shortest Euclidean paths in raw pixel space.
\end{itemize}
The model is an optimized CNN with convolutional layers (64 and 32 filters) and a 128-dimensional bottleneck.

\paragraph{NLP Benchmarks:}
We utilize IMDB (binary reviews), SST-2 (binary), and SST-5 (5-class sentiment). Reviews are represented as 50-dimensional average-pooled Bag-of-Words (BoW) embeddings. We utilize human-edited counterfactuals for IMDB \cite{Kaushik2020Learning} and automated counterfactuals for the SST tasks. Scarcity was simulated with subsets ranging from 10 to 200 samples.

\subsection{Active Learning Integration}
To evaluate the impact of sample quality, we compared two data acquisition strategies: (1) \textbf{AL ON}: Selecting boundary samples via entropy-based importance sampling with a diversity parameter; (2) \textbf{AL OFF}: Purely random selection. This comparison allows us to isolate the synergy between informative data selection and geometric supervision.

\section{Results and Analysis}
\label{sec:results}

\subsection{NLP Benchmarks and AL Synergy}
As shown in Table \ref{tab:nlp_results}, CGA achieves substantial gains on IMDB and SST-5 sentiment tasks, consistently outperforming both standard Cross Entropy and the GS baseline. The Softplus variant was the most robust performer on the high-dimensional IMDB dataset, achieving a peak accuracy of 63.11\% with only 200 samples, representing a +7.17\% gain over the baseline.

\begin{table}[h]
\centering
\caption{Validation Accuracy on NLP benchmarks (AL ON, 5 Epochs). CGA consistently outperforms Gradient Supervision (GS).}
\label{tab:nlp_results}
\begin{tabular}{@{}lccccc@{}}
\toprule
\textbf{Dataset} & \textbf{Size} & \textbf{Baseline} & \textbf{GS \cite{teney2020learning}} & \textbf{Best CGA} & \textbf{Gain} \\ \midrule
IMDB & 200 & 55.94\% & 60.82\% & \textbf{63.11\% (SP)} & \textbf{+7.17\%}
 \\
IMDB & 100 & 51.43\% & 56.73\% & \textbf{56.56\% (Sign)} & +5.13\%
 \\
SST-5 & 100 & 53.06\% & 56.46\% & \textbf{56.46\% (ReLU)} & +3.40\%
 \bottomrule
\end{tabular}
\end{table}

However, the utility of CGA was found to be absolutely dependent on sample quality. As shown in Table \ref{tab:al_synergy}, disabling Active Learning caused both CGA and GS to perform worse than the baseline on text benchmarks. This suggests that directional intuition on redundant samples introduces noise that conflicts with the primary classification signal. CGA utility actually flips from positive to negative without AL.

\begin{table}[h]
\centering
\caption{Synergy between Alignment and Active Learning (Size 100). CGA utility flips from positive to negative on random samples.}
\label{tab:al_synergy}
\begin{tabular}{lcccc}
\toprule
\textbf{Dataset} & \textbf{AL ON (CGA)} & \textbf{AL ON (GS)} & \textbf{AL OFF (CGA)} & \textbf{AL OFF (GS)} \\ \midrule
IMDB & \textbf{+5.13\%} & +5.30\% & -2.26\% & -3.12\%
 \\
SST-2 & \textbf{+1.36\%} & -1.36\% & -2.72\% & -3.50\%
 \\
SST-5 & \textbf{+3.40\%} & +3.40\% & -4.77\% & -5.20\%
 \bottomrule
\end{tabular}
\end{table}

\subsection{Complex Topologies: Synthetic 2D Benchmarks}
CGA's ability to regularize complex topologies was most evident in the Concentric Circles task. Standard Cross Entropy struggle to define the circular boundary from sparse points, often collapsing to a biased linear separation. In contrast, CGA (Sign variant) reached 72.00\% accuracy, a +14.50\% improvement. Interestingly, the standard GS method struggled with these synthetic boundaries in low-data regimes, likely due to high sensitivity of angular alignment to local gradient noise (Table \ref{tab:2d_results}).

\begin{table}[h]
\centering
\caption{Validation Accuracy on Synthetic 2D Topologies. CGA variants provide much higher stability than standard GS.}
\label{tab:2d_results}
\begin{tabular}{@{}lccccc@{}}
\toprule
\textbf{Dataset} & \textbf{Size} & \textbf{Baseline (CE)} & \textbf{GS \cite{teney2020learning}} & \textbf{Best CGA} & \textbf{Gain} \\ \midrule
Circles & 100 & 57.50\% & 39.00\% & \textbf{72.00\% (Sign)} & \textbf{+14.50\%}
 \\
XOR     & 30 & \textbf{97.50\%} & 70.00\% & 93.00\% (SP) & -4.50\%
 \\
TwoMoons & 30 & \textbf{82.00\%} & 63.00\% & 82.00\% (Tie) & 0.00\%
 \bottomrule
\end{tabular}
\end{table}

\subsection{Path Feasibility Analysis on MNIST}
We evaluated how the feasibility of the counterfactual path affects alignment on MNIST. FACE paths, which prioritize plausible data transitions by constructing density-weighted graphs, provided superior results at larger training sizes compared to raw geometric paths. This confirms that respecting the underlying data manifold when providing expert directions is crucial for high-dimensional inputs. We also included PCA-based paths as a baseline, which performed similarly to raw paths but slightly better in some cases (Table \ref{tab:mnist_paths}).

\begin{table}[h]
\centering
\caption{MNIST Accuracy Comparison across Path Types (AL ON, Size 200). FACE paths provide the highest utility for CGA.}
\label{tab:mnist_paths}
\begin{tabular}{lcccc}
\toprule
\textbf{Path Type} & \textbf{Baseline (CE)} & \textbf{GS \cite{teney2020learning}} & \textbf{CGA (Ours)} & \textbf{Gain} \\ \midrule
Raw & 83.50\% & 83.50\% & 84.52\% (Sign) & +1.02\%
 \\
PCA & 83.50\% & 83.60\% & 84.56\% (SP) & +1.06\%
 \\
FACE & 83.50\% & 83.72\% & \textbf{84.78\% (SP)} & \textbf{+1.28\%}
 \bottomrule
\end{tabular}
\end{table}

\section{Discussion}
\label{sec:discussion}

\subsection{Smooth Alignment and Margin Consistency}
The consistent superiority of the \textbf{Softplus} and \textbf{Sign} CGA variants over both standard CE and Gradient Supervision suggests that explicitly penalizing the directional derivative is a more robust signal than pure angular alignment. While GS aligns the gradient direction, it often fails to encourage a definitive drop in class confidence, which is critical for forming sharp decision boundaries from sparse data. Softplus ensures the model maintains a high-margin drop toward the counterfactual boundary even when "correct," leading to more stable generalization.

\subsection{Active Learning as a Geometric Filter}
The failure of CGA on random data subsets (Table \ref{tab:al_synergy}) reinforces the view that expert intuition is a high-precision, low-recall signal. It provides the model with structural information that categorical labels cannot convey, but only when applied to the most ambiguous regions of the feature space. In central regions, directional intuition can conflict with the primary classification signal. The synergy between Active Learning and CGA is therefore fundamental to its efficacy.

\subsection{Expert Effort and Annotation Budget}
While CGA increases the annotation burden per sample (providing a direction vs just a label), it significantly reduces the number of samples required to reach a target accuracy. The results show that a model can inherit complex structural priors from a few expert-annotated boundary cases. This provides a blueprint for human-AI teaching loops that optimize for total expert effort rather than just label count.

\section{Conclusions}
We have presented Counterfactual Gradient Alignment (CGA), a framework for supervising model gradients via expert directional intuition. Our evaluation across diverse benchmarks proves that smooth directional losses combined with Active Learning significantly improve generalization in data-scarce environments. By moving beyond categorical labels toward machine-interpretable geometric supervision, we provide a path toward more collaborative and data-efficient learning. Future research will investigate the application of CGA to foundation models and multi-modal grounding.

\balance
\bibliographystyle{ACM-Reference-Format}
\bibliography{sample}

\appendix
\section{Importance Sampling Implementation}
To ensure that our framework selects the most informative samples for counterfactual annotation while maintaining diversity across the feature space, we employ an entropy-based importance sampling strategy with an integrated diversity parameter. This approach addresses the common limitations of purely uncertainty-based or diversity-based sampling methods by combining their respective strengths \cite{he2014active,doucet2024bridging}.

\textbf{Implementation Details.} Our sampling procedure operates iteratively during training. At each selection round, we first compute the softmax probabilities from the current model's logits: $\mathbf{p}_i = \text{softmax}(\mathbf{f}_\theta(\mathbf{x}_i))$. The uncertainty of each sample is then quantified using the entropy of these probability distributions: $U(\mathbf{x}_i) = -\sum_{y \in \mathcal{Y}} p_{i,y} \log p_{i,y}$, where $p_{i,y}$ denotes the predicted probability for class $y$. Samples with high entropy indicate regions of significant uncertainty.

To address redundancy, we incorporate a diversity parameter $\beta$ that encourages the selection of samples spanning different regions. We use the learned embeddings $\mathbf{e}_i$ from our model's intermediate representations to compute representativeness $Rep(\mathbf{x}_i)$ based on similarity to other unlabeled instances: $Rep(\mathbf{x}_i) = \frac{1}{|U|} \sum_{\mathbf{x}_j \in U} \exp\left(-\frac{||\mathbf{e}_i - \mathbf{e}_j||^2}{2\sigma^2}\right)$. The combined information content is then: $Info(\mathbf{x}_i) = U(\mathbf{x}_i)^\alpha \cdot Rep(\mathbf{x}_i)^\beta$.

\section{Shortest path definitions}
\paragraph{Raw Distance Counterfactuals}
Generates counterfactuals by finding the nearest instance in the target class using standard Euclidean distance: $\mathbf{x}_{cf} = \arg\min_{\mathbf{x} \in \mathcal{X}_{target}} ||\mathbf{x} - \mathbf{x}_0||_2$. This approach provides the most direct path but may traverse low-density regions.

\paragraph{PCA Distance Counterfactuals}
Addresses feature correlation by transforming data into principal component space before finding the nearest counterfactual: $\mathbf{z}_{cf} = \arg\min_{\mathbf{z} \in \mathcal{Z}_{target}} ||\mathbf{z} - \mathbf{z}_0||_2$, where $\mathbf{z} = \mathbf{P}^T\mathbf{x}$. This can reveal semantically meaningful proximity relationships obscured by raw feature scales.

\end{document}
%%%%%%%%%%%%%%%%%%%%%%%%%%%%%%%%%%%%%%%%%%%%%%%%%%%%%%%%%%%%%%%%%%%%%%%%